\chapter{绪论}
\section{研究背景与意义}
高校实验室通常同时管理仪器、耗材和教学套件。设备借用涉及申请、审批、领取、归还和状态复核等环节，任何一个环节缺少记录，都可能造成设备去向不清或库存数据滞后。小规模实验室可以通过纸质登记维持运行，但当设备种类和使用人数增加后，人工汇总会带来重复录入、状态更新不及时和历史记录难以检索等问题。

设备借用系统的价值不只是把纸质表格搬到网页上，更在于明确业务状态和操作责任。系统需要在提交申请时校验可借数量，在审批时保留处理人，在归还时记录设备状态，并把关键操作写入审计日志。这样，管理员能够从一条申请追踪到完整的设备流转过程，学生也能清楚了解申请处于何种状态。
\section{相关研究与实践}
信息管理系统通常将权限控制、业务状态和数据审计作为基础能力。面向资产管理的研究强调统一编码和生命周期记录，面向流程管理的研究则关注状态转换、并发处理和异常回滚。本文结合实验室借用场景，将设备台账与借用流程放在同一数据模型中，并通过事务处理降低库存状态不一致的风险\cite{ref2}。
\section{研究内容与论文结构}
本文围绕“登记、申请、审批、归还、审计”五个环节展开。第一章介绍研究背景和工作内容；第二章分析业务需求；第三章说明系统架构、数据模型和流程设计；第四章给出主要模块的详细实现；第五章通过测试用例分析系统行为；第六章总结工作并讨论改进方向。

